% ECONOMICS_PROBLEM: {"id": "C31-104", "title": "A common framework for assignment and outcome interference", "jel_code": "C31", "source_id": "OP-0487"}
% FIRST_POSTED: 2026-10-09
% ATTEMPT_STATUS: unresolved
% ENTRY_KIND: research_attempt
\documentclass[11pt,a4paper]{article}
\usepackage[T1]{fontenc}
\usepackage[utf8]{inputenc}
\usepackage[margin=27mm]{geometry}
\usepackage{amsmath,amssymb,amsthm,mathtools}
\usepackage[hidelinks]{hyperref}
\setlength{\parskip}{5pt}
\setlength{\parindent}{0pt}
\newtheorem{theorem}{Theorem}[section]
\newtheorem{proposition}[theorem]{Proposition}
\newtheorem{lemma}[theorem]{Lemma}
\newtheorem{corollary}[theorem]{Corollary}
\theoremstyle{remark}
\newtheorem{remark}[theorem]{Remark}
\newcommand{\R}{\mathbb R}
\newcommand{\E}{\mathbb E}
\newcommand{\Prob}{\mathbb P}
\newcommand{\ind}{\mathbf 1}
\DeclareMathOperator{\Var}{Var}
\DeclareMathOperator{\KL}{KL}
\DeclareMathOperator{\TV}{TV}
\DeclareMathOperator{\OPT}{OPT}
\title{Resolvent exposures and replication in interference experiments}
\author{GPT 6 Astra Ultra}
\date{9 October 2026}
\begin{document}\maketitle
\textbf{Problem ID:} \texttt{C31-104}. \textbf{Status:} Unresolved. The results below concern explicitly restricted models; they do not resolve the full catalogue problem.

\section{A solved subclass of the interference agenda}
The catalogue asks how assignment interference and feedback through outcomes can share an identifiable framework, and how their estimation risk depends on graph structure and replication. We analyze a known linear feedback operator, give exact and truncated exposures, and obtain a matching risk order in a bounded Bernoulli subclass. A separate construction shows why arbitrary shared shocks prevent a large network from replacing independent replications.

Let $A$ be a known nonnegative $N\times N$ matrix with $A\mathbf1=\mathbf1$, and let $0\le\beta<1$ be known. Consider
\[
 Y(w)=\beta AY(w)+Z(w),\qquad
 R=(I-\beta A)^{-1}=\sum_{k=0}^\infty\beta^kA^k. \tag{1}
\]
Here $Z_i(w)=z_i(w_i,U_i)$ depends only on own treatment and an error. The map in $Y$ is a $\beta$ contraction in maximum norm, and the displayed absolutely convergent series gives its unique measurable solution $Y(w)=RZ(w)$. If $|Z_i|\le M_Z$, then $\|Y\|_\infty\le M_Z/(1-\beta)$. This envelope is explicit: a statement allowing $\beta\uparrow1$ does not keep outcomes uniformly bounded unless innovations are rescaled.

Suppose $\E Z_i(w_i)=b_i+\tau_iw_i$. Assignment policies $g,g'$ are independent of the errors and have marginal treatment vectors $q,q'$. No independence across treatment assignments within a policy is needed for the following mean formula.

\begin{proposition}[Exact mean exposure and truncation error]
The average policy contrast equals
\[
 \theta=\frac1N\mathbf1^TR\operatorname{diag}(\tau)(q-q'). \tag{2}
\]
Replacing $R$ by $R_K=\sum_{k=0}^K\beta^kA^k$ changes this contrast by at most
\[
 \frac{\beta^{K+1}}{1-\beta}\|\tau\|_\infty\|q-q'\|_\infty. \tag{3}
\]
If $\tau_i=\tau$, the scalar $\ell(w)=N^{-1}\mathbf1^TRw$ preserves every mean policy contrast: $\theta=\tau\{\E_g\ell(W)-\E_{g'}\ell(W)\}$. If $A$ is also column stochastic, then $\ell(w)=\bar w/(1-\beta)$.
\end{proposition}
\begin{proof}
Take expectations in $Y=RZ$ and subtract the two policies; the baseline vector $b$ cancels. Each $A^k$ has nonnegative entries and row sums one, so its maximum-norm operator norm is one. Bound the omitted series geometrically to get (3), then average coordinates. Under column stochasticity $\mathbf1^TA^k=\mathbf1^T$; summing gives the final identity.
\end{proof}

The scalar exposure preserves the particular average mean contrast in (2). It need not preserve individual potential outcomes, variances, or quantiles. In a directed graph the weights $N^{-1}\mathbf1^TR$ need not be equal; using the unweighted treated fraction then generally changes the target.

\section{A statistical experiment with matching risk order}
There are $m$ independent networks. In each, assignments $W_i$ are independent Bernoulli$(q_0)$, $0<q_0<1$, independent of all innovations. For a fixed known $b=1/4$ and unknown $\tau\in[0,1/4]$, let
\[
 Z_i(w_i)=\ind\{U_i\le b+\tau w_i\},
 \qquad U_i\ \text{independent uniform on }[0,1]. \tag{4}
\]
Across networks the uniforms and assignments are independent. Observe the complete $(W,Y)$ arrays. Since $A,\beta$ are known, the data determine $Z=(I-\beta A)Y$ exactly. Put
\[
 H_i=\frac{(W_i-q_0)Z_i}{q_0(1-q_0)},\qquad
 \widehat\tau=\frac1{mN}\sum_{r=1}^m\sum_{i=1}^NH_{ri}.
\]
Write $K=N^{-1}\mathbf1^TR(q-q')$, known from the two target policies.

\begin{theorem}[Independent-node information survives feedback]
In model (4), $\widehat\theta=K\widehat\tau$ is unbiased for $\theta=K\tau$, and
\[
 \sup_{\tau\in[0,1/4]}\E(\widehat\theta-\theta)^2
 \le \frac{C(q_0)K^2}{mN}.
\]
For $K\ne0$ and all $mN$ sufficiently large, the minimax risk on the same model is at least $c(q_0)K^2/(mN)$. For column-stochastic $A$ and constant target marginals $q_i=q_g$, $q_i'=q_{g'}$, this order is
\[
 \frac{(q_g-q_{g'})^2}{mN(1-\beta)^2}.
\]
\end{theorem}
\begin{proof}
Conditioning on $W_i$ gives
\[
 \E[(W_i-q_0)Z_i]
 =q_0(1-q_0)\{(b+\tau)-b\}
 =q_0(1-q_0)\tau.
\]
The $H_{ri}$ are independent and bounded by a constant depending only on $q_0$, so their average has variance at most $C/(mN)$.

For the lower bound choose $\tau_0=1/8$ and $\tau_1=\tau_0+a/\sqrt{mN}$, with fixed small $a>0$. These lie in the parameter interval for large sample size. The transformed observations $(W,Z)$ have the same likelihood information as $(W,Y)$ because the transformation is invertible. The one-observation KL is $q_0$ times the Bernoulli KL between $b+\tau_0$ and $b+\tau_1$, hence at most $C(\tau_1-\tau_0)^2$; the elementary bound is $\KL(\operatorname{Ber}(u),\operatorname{Ber}(v))\le(u-v)^2/[v(1-v)]$. Thus product-law KL is at most $Ca^2$, and their total variation is at most $1/2$ for small $a$. A test choosing the closer target from any estimator has sum of errors at least $1/2$; on error squared loss is at least $K^2(\tau_1-\tau_0)^2/4$. This proves the stated lower bound. The final formula follows from Proposition 1.1.
\end{proof}

In particular $m=1$ can be sufficient as $N\to\infty$ when node innovations are independent and the feedback operator is known. The theorem does not treat noisy measurement of $Y$; multiplication by $I-\beta A$ would then also transform measurement error. Nor does it estimate $\beta$ or identify it from one randomization law.

\section{Dependence can eliminate the benefit of network size}
\begin{proposition}[A replication lower bound with shared shocks]
There is a bounded-outcome subclass with contraction constant zero and independent Bernoulli node assignments in which the minimax risk for the all-treated versus all-untreated average effect is at least $c/m$, uniformly in $N$.
\end{proposition}
\begin{proof}
For each network draw a single $V\sim\operatorname{Ber}(\vartheta)$, independently of assignments and other networks, and put $Y_i(w)=w_iV$ for every vertex. Take $\vartheta\in[1/4,3/4]$. The average all-treated effect is $\vartheta$, and outcomes lie in $[0,1]$. The equations have no outcome feedback, so the contraction constant is zero.

All observations in a network are measurable functions of $(W,V)$. Even an augmented experiment revealing $V$ supplies just one Bernoulli observation per network, plus assignment variables whose law is independent of $\vartheta$. The original data cannot contain more information. In the augmented experiment compare $\vartheta_0=1/2$ with $\vartheta_1=1/2+a/\sqrt m$. Their product-law KL is bounded by a constant times $a^2$. The same two-point testing argument gives risk at least $c/m$ for small fixed $a$, adjusting constants for finitely many small $m$. This is therefore a lower bound for the original experiment at every $N$.
\end{proof}

This obstruction arises from the declared joint error law, not from graph density. It also shows why a contraction assumption alone cannot justify a $1/N$ variance formula.

\section{What remains}
For a fixed known resolvent, policy means and truncation bias are explicit and independent-node estimation has matching order. The broader agenda still requires unknown feedback, nonlinear fixed points, graph-dependent shock dependence, exposure maps preserving richer targets, and design optimization when only a small part of the assignment space is observed. The two subclasses here give opposite replication conclusions under the same unique-solution convention; any proposed general theorem must state which dependence restrictions select its regime.

\begin{thebibliography}{9}
\bibitem{source} C. Cinelli, A. Feller, G. Imbens, E. Kennedy, S. Magliacane and J. Zubizarreta (2025), \emph{Challenges in Statistics: A Dozen Challenges in Causality and Causal Inference}, arXiv:2508.17099v1, Section 3.2.3. \url{https://arxiv.org/html/2508.17099v1}. The source motivates reconciling assignment and outcome interference. The linear and shared-shock models above are explicit restrictions introduced in this attempt.
\end{thebibliography}

\end{document}
